/home/sybil/.tex_templates/template.tex
\setupHW{MATH 341: Howework 5}
\begin{document}
\maketitle
\section{Question 1}
Let $A$ be defined as:
\[
	A :=
	\begin{bmatrix}
		4 & 0 & 1 \\
		0 & 2 & 0 \\
		1 & 0 & 4
	\end{bmatrix}
\]
This matrix is symmetric. As a hint, $\lambda=2$ is one eigenvalue.
\begin{enumerate}
	\item[(a)] Find all eigenvalues of $A$.
	\item[(b)] Find an eigenvector for each eigenvalue.
	\item[(c)] Verify that your eigenvectors are mutually orthogonal.
	\item[(d)] Normalize your eigenvectors to form an orthogonal matrix $Q$, and write down a diagonal matrix $\Lambda$ such that \[ A = Q\Lambda Q^T. \]
\end{enumerate}
\subsection{Find all eigenvalues of $A$.}
\[
	A-\lambda I :=
	\begin{bmatrix}
		4-\lambda & 0         & 1         \\
		0         & 2-\lambda & 0         \\
		1         & 0         & 4-\lambda
	\end{bmatrix}
\]
Determinant:
\[
	(2-\lambda) \det \begin{bmatrix}
		4-\lambda & 1         \\
		1         & 4-\lambda
	\end{bmatrix}
\]
Characteristic Equation:
\[
	\begin{aligned}
		0          & = (2-\lambda) \left[ (4-\lambda)^2 - 1 \right]              \\
		0          & = (2-\lambda)(4-\lambda)^2 - (2-\lambda)                    \\
		0          & = (2-\lambda)(16 -8\lambda + \lambda^2) - (2-\lambda)       \\
		0          & = (2-\lambda)[(16 -8\lambda + \lambda^2) - 1]               \\
		0          & = (2-\lambda)[15 -8\lambda + \lambda^2 ]                    \\
		0          & = (2-\lambda)(\lambda-3)(\lambda-5)                         \\
		\therefore & \quad \lambda_2 = 2 \quad \lambda_3 = 3 \quad \lambda_5 = 5
	\end{aligned}
\]
\subsection{Find an eigenvector for each eigenvalue.}
\begin{center}
	Let $M_n := A-\lambda_n I ~.~ n \in \{ 2,3,5 \}$ \quad\quad Let $V_n := \mathscr{B}[Nul(M_n)] ~.~ n \in \{ 2,3,5 \}$
\end{center}
\[
	A-\lambda_2 I =
	\begin{bmatrix}
		2 & 0 & 1 \\
		0 & 0 & 0 \\
		1 & 0 & 2
	\end{bmatrix} \implies
	Nul(A-\lambda_2 I) = \begin{bmatrix}1&0&0\\ 0&0&1\\ 0&0&0\end{bmatrix}
	\implies
	\begin{cases}
		x_2 & = 0 \\
		x_2 & = 1 \\
		x_2 & = 0
	\end{cases}
	\implies V_1 = \left\{ \begin{bmatrix} 0 \\ 1 \\0 \end{bmatrix} \right\}
\]

\[
	A-\lambda_3 I :=
	\begin{bmatrix}
		1 & 0  & 1 \\
		0 & -1 & 0 \\
		1 & 0  & 1
	\end{bmatrix} \implies
	Nul(A-\lambda_3 I) = \begin{bmatrix}1&0&1\\ 0&1&0\\ 0&0&0\end{bmatrix}
	\implies \begin{cases}
		x_3 & = -1 \\
		x_3 & = 0  \\
		x_3 & = 1
	\end{cases}
	\implies V_3 = \left\{ \begin{bmatrix} -1 \\ 0 \\1 \end{bmatrix} \right\}
\]

\[
	A-\lambda_5 I :=
	\begin{bmatrix}
		-1 & 0  & 1  \\
		0  & -3 & 0  \\
		1  & 0  & -1
	\end{bmatrix}
	Nul(A-\lambda_5 I) = \begin{bmatrix}1&0&-1\\ 0&1&0\\ 0&0&0\end{bmatrix}
	\implies \begin{cases}
		x_3 & = 1 \\
		x_3 & = 0 \\
		x_3 & = 1
	\end{cases}
	\implies V_5 = \left\{ \begin{bmatrix} 1 \\ 0 \\1 \end{bmatrix} \right\}
\]

\[ V_1 = \begin{bmatrix} 0 \\ 1 \\ 0 \end{bmatrix} V_3 = \begin{bmatrix} -1 \\ 0 \\1 \end{bmatrix} V_5 = \begin{bmatrix} 1 \\ 0 \\1 \end{bmatrix} \]

\subsection{Verify that your eigenvectors are mutually orthogonal.}
\[
	V_1^TV_3 :=
	\begin{bmatrix} 0 & 1 & 0 \end{bmatrix}
	\begin{bmatrix} -1 \\ 0 \\1 \end{bmatrix} = 0
	\quad \quad
	V_1^TV_5
	\begin{bmatrix} 0 & 1 & 0 \end{bmatrix}
	\begin{bmatrix} 1 \\ 0 \\1 \end{bmatrix} = 0
	\quad \quad
	V_3^TV_5
	\begin{bmatrix} -1 & 0 & 1 \end{bmatrix}
	\begin{bmatrix} 1 \\ 0 \\1 \end{bmatrix} = 0
\]


\subsection{Normalize your eigenvectors to form an orthogonal matrix $Q$, and write down a diagonal matrix $\Lambda$ such that $A = Q\Lambda Q^T$}
\[
	V_1 = \begin{bmatrix} 0 \\ 1 \\ 0 \end{bmatrix} \implies q_1 = \begin{bmatrix} 0 \\ 1 \\ 0 \end{bmatrix} \quad\quad\quad
	V_3 = \begin{bmatrix} -1 \\ 0 \\1 \end{bmatrix} \implies q_3 = \begin{bmatrix} -\frac{1}{\sqrt{2}} \\ 0 \\\frac{1}{\sqrt{2}} \end{bmatrix} \quad\quad\quad
	V_5 = \begin{bmatrix} 1 \\ 0 \\1 \end{bmatrix}  \implies q_5 = \begin{bmatrix} \frac{1}{\sqrt{2}} \\ 0 \\\frac{1}{\sqrt{2}} \end{bmatrix}
\]

\[
	A = \begin{bmatrix}
		4 & 0 & 1 \\
		0 & 2 & 0 \\
		1 & 0 & 4
	\end{bmatrix} \quad\quad
	Q =
	\begin{bmatrix}
		0 & -\frac{1}{\sqrt{2}} & \frac{1}{\sqrt{2}} \\
		1 & 0                   & 0                  \\
		0 & \frac{1}{\sqrt{2}}  & \frac{1}{\sqrt{2}}
	\end{bmatrix} \quad\quad
	\Lambda = \begin{bmatrix}
		2 & 0 & 0 \\
		0 & 3 & 0 \\
		0 & 0 & 5
	\end{bmatrix} \quad\quad
	Q^T =
	\begin{bmatrix}
		0                   & 1 & 0                  \\
		-\frac{1}{\sqrt{2}} & 0 & \frac{1}{\sqrt{2}} \\
		\frac{1}{\sqrt{2}}  & 0 & \frac{1}{\sqrt{2}}
	\end{bmatrix}
\]

Final expression:

\[
	\begin{bmatrix}
		4 & 0 & 1 \\
		0 & 2 & 0 \\
		1 & 0 & 4
	\end{bmatrix} =
	\begin{bmatrix}
		0 & -\frac{1}{\sqrt{2}} & \frac{1}{\sqrt{2}} \\
		1 & 0                   & 0                  \\
		0 & \frac{1}{\sqrt{2}}  & \frac{1}{\sqrt{2}}
	\end{bmatrix}
	\begin{bmatrix}
		2 & 0 & 0 \\
		0 & 3 & 0 \\
		0 & 0 & 5
	\end{bmatrix}
	\begin{bmatrix}
		0                   & 1 & 0                  \\
		-\frac{1}{\sqrt{2}} & 0 & \frac{1}{\sqrt{2}} \\
		\frac{1}{\sqrt{2}}  & 0 & \frac{1}{\sqrt{2}}
	\end{bmatrix}
\]

\section{Question 2}
Let $A$ be defined as:
\[
	A = \begin{bmatrix}
		1 & 1 & 0 \\
		0 & 1 & 0 \\
		0 & 0 & 2
	\end{bmatrix}
\]
\begin{enumerate}
	\item[(a)] Find all eigenvalues of $A$ and their algebraic multiplicities.
	\item[(b)] Find the eigenspace corresponding to each eigenvalue.
	\item[(c)] Find all eigenpairs of $A$.
	\item[(d)] Decide whether $A$ is diagonalizable. Explain briefly.
\end{enumerate}
\[
	A-\lambda I =
	\begin{bmatrix}
		1 & 1 & 0 \\
		0 & 1 & 0 \\
		0 & 0 & 2
	\end{bmatrix}
\]

\subsection{Find all eigenvalues of $A$ and their algebraic multiplicities.}
\[
	A-\lambda I = \begin{bmatrix}
		1-\lambda & 1         & 0         \\
		0         & 1-\lambda & 0         \\
		0         & 0         & 2-\lambda
	\end{bmatrix}
\]
\[
	\det(A - \lambda I) = \begin{vmatrix}
		1-\lambda & 1         & 0         \\
		0         & 1-\lambda & 0         \\
		0         & 0         & 2-\lambda
	\end{vmatrix}
\]
Determinant:
\[
	\begin{aligned}
		0 = \det(A - \lambda I) & = (1-\lambda) \begin{vmatrix} 1-\lambda & 0 \\ 0 & 2-\lambda \end{vmatrix} - 1 \begin{vmatrix} 0 & 0 \\ 0 & 2-\lambda \end{vmatrix} + 0 \begin{vmatrix} 0 & 1-\lambda \\ 0 & 0 \end{vmatrix} \\
		0                       & = (1-\lambda)[(1-\lambda)(2-\lambda) - 0] - 0                                                                                                                                                \\
		0                       & = (1-\lambda)^2(2-\lambda)
	\end{aligned}
\]
Characteristic Equation:
\[ (1-\lambda)^2 (2-\lambda) = 0 \]
\[
	\begin{aligned}
		\implies \lambda_1 & = 1 \quad \text{(Algebraic Multiplicity = 2)} \\
		\implies \lambda_2 & = 2 \quad \text{(Algebraic Multiplicity = 1)}
	\end{aligned}
\]

\subsection{Find the eigenspace corresponding to each eigenvalue.}
\[
	A-\lambda I = \begin{bmatrix}
		1-\lambda & 1         & 0         \\
		0         & 1-\lambda & 0         \\
		0         & 0         & 2-\lambda
	\end{bmatrix}
\]
\[
	\begin{aligned}
		\implies \lambda_1 & = 1 \quad \text{(Algebraic Multiplicity = 2)} \\
		\implies \lambda_2 & = 2 \quad \text{(Algebraic Multiplicity = 1)}
	\end{aligned}
\]

\subsection{Find the eigenspace corresponding to each eigenvalue.}
\begin{center}
	Let $M_n := A-\lambda_n I ~.~ n \in \{ 2,3,5 \}$ \quad\quad Let $V_n := \mathscr{B}[Nul(M_n)] ~.~ n \in \{ 2,3,5 \}$
\end{center}
\[
	A-\lambda I = \begin{bmatrix}
		1-\lambda & 1         & 0         \\
		0         & 1-\lambda & 0         \\
		0         & 0         & 2-\lambda
	\end{bmatrix}
\]

\[
	A-\lambda_1 I =
	\begin{bmatrix}
		0 & 1 & 0 \\
		0 & 0 & 0 \\
		0 & 0 & 1
	\end{bmatrix}
	Nul(A-\lambda_1 I) =
	\begin{bmatrix}
		0 & 1 & 0 \\
		0 & 0 & 0 \\
		0 & 0 & 1
	\end{bmatrix} \implies
	\begin{cases}
		x_1 = 1 \\
		x_1 = 0 \\
		x_1 = 0 \\
	\end{cases} \implies
	V_1 = \left\{
	\begin{bmatrix} 1 \\ 0 \\ 0 \end{bmatrix}
	\right\}
\]

\[
	A-\lambda_2 I =
	\begin{bmatrix}
		-1 & 1  & 0 \\
		0  & -1 & 0 \\
		0  & 0  & 0
	\end{bmatrix} \implies Nul(A-\lambda_2 I) =
	\begin{bmatrix}1&0&0\\ 0&1&0\\ 0&0&0\end{bmatrix} \implies
	\begin{cases}
		x_3 = 0 \\
		x_3 = 0 \\
		x_3 = 1 \\
	\end{cases} \implies
	V_2 = \left\{
	\begin{bmatrix} 0 \\ 0 \\ 1 \end{bmatrix}
	\right\}
\]
\subsection{Find all eigenpairs of $A$.}
\[
	(\lambda_1, V_1) \quad (\lambda_2, V_2)
\]
\subsection{Decide whether $A$ is diagonalizable. Explain briefly.}
\[
	\text{Geometric Multiplicity} < \text{Algebraic Multiplicity} \implies A \text{ is NOT Diagonalizable}
\]

\section{Question 3}
Let $A$ be defined as:
\[ A =
	\begin{bmatrix}
		1 & -2 \\
		2 & 1
	\end{bmatrix}
\]
\begin{enumerate}
	\item[(a)] Find the characteristic polynomial of $A$.
	\item[(b)] Find the eigenvalues of $A$ over $\mathbb{C}$.
	\item[(c)] For each eigenvalue, find an eigenvector in $\mathbb{C}^2$.
	\item[(d)] Explain why $A$ has no real eigenvectors.
\end{enumerate}
\subsection{Find the characteristic polynomial of $A$.}
\[ A-\lambda I =
	\begin{bmatrix}
		1-\lambda & -2        \\
		2         & 1-\lambda
	\end{bmatrix}
\]

Determinant:
\[
	\det(A-\lambda I) =
	\begin{vmatrix}
		1-\lambda & -2        \\
		2         & 1-\lambda
	\end{vmatrix}
	= (1-\lambda)(1-\lambda) - (-2)\cdot2 = (1-\lambda)^2 + 4
\]
\subsection{Find the eigenvalues of $A$ over $\mathbb{C}$.}
\[
	\begin{align}
		(1-\lambda)^2 + 4            & = 0 \\
		1 - 2\lambda + \lambda^2 + 4 & = 0 \\
		5 - 2\lambda + \lambda^2     & = 0 \\
	\end{align}
\]
Therefore there must exists roots $\in \mathbb{C}$. Obtain complex conjugate:
\[
	\text{Quadratic Formula:} \quad \lambda & = \frac{2 \pm \sqrt{(-2)^2 - 4 \cdot 1 \cdot 5}}{2 \cdot 1} = 1 \pm 2i \\
\]
Thus we obtain the following:
\[ \lambda_1 = 1+2i \quad \lambda_2 = 1-2i \]
\subsection{For each eigenvalue, find an eigenvector in $\mathbb{C}^2$.}
\begin{center}
	Let $M_n := A-\lambda_n I ~.~ n \in \{ 2,3,5 \}$ \quad\quad Let $V_n := \mathscr{B}[Nul(M_n)] ~.~ n \in \{ 2,3,5 \}$
\end{center}

For $\lambda_1 = 1+2i$:
\[
	M_1 = A-(1+2i)I =
	\begin{bmatrix}
		1-(1+2i) & -2       \\
		2        & 1-(1+2i)
	\end{bmatrix} =
	\begin{bmatrix}
		-2i & -2  \\
		2   & -2i
	\end{bmatrix}
\]
To find the null space, we solve $M_1 \mathbf{v} = 0$:
\[
	-2ix - 2y = 0 \implies 2y = -2ix \implies y = -ix
\]
Choosing $x=1$, we get $y=-i$.
\[
	V_1 = \left\{ \begin{bmatrix} 1 \\ -i \end{bmatrix} \right\}
\]

For $\lambda_2 = 1-2i$:
\[
	M_2 = A-(1-2i)I =
	\begin{bmatrix}
		1-(1-2i) & -2       \\
		2        & 1-(1-2i)
	\end{bmatrix} =
	\begin{bmatrix}
		2i & -2 \\
		2  & 2i
	\end{bmatrix}
\]
To find the null space, we solve $M_2 \mathbf{v} = 0$:
\[
	2ix - 2y = 0 \implies 2y = 2ix \implies y = ix
\]
Choosing $x=1$, we get $y=i$.
\[
	V_2 = \left\{ \begin{bmatrix} 1 \\ i \end{bmatrix} \right\}
\]

\subsection{Explain why $A$ has no real eigenvectors.}
A matrix $A$ has a real eigenvector $\mathbf{v} \in \mathbb{R}^2$ if and only if there exists a real eigenvalue $\lambda \in \mathbb{R}$ such that $A\mathbf{v} = \lambda\mathbf{v}$.

If we assume $\mathbf{v}$ is a real, non-zero vector, then the product $A\mathbf{v}$ must result in another real vector. For the equation $A\mathbf{v} = \lambda\mathbf{v}$ to hold, $\lambda$ would necessarily have to be a real number. Because the eigenvalues of this matrix are strictly complex ($\lambda = 1 \pm 2i$), there is no real scalar $\lambda$ that can satisfy the eigenvalue equation for any non-zero real vector $\mathbf{v}$.

Therefore, any eigenvector of $A$ must necessarily have complex components, placing it in $\mathbb{C}^2 \setminus \mathbb{R}^2$.

\section{Question 4}
Find one singular value decomposition of the $3 \times 2$ matrix. Let $A$ be defined as:
\[ A = \begin{bmatrix} 1 & 1 \\ 1 & 1 \\ 0 & 0 \end{bmatrix}. \]
More precisely:
\begin{enumerate}
	\item[(a)] Compute $A^T A$ and find its eigenvalues.
	\item[(b)] Find the singular values of $A$.
	\item[(c)] Find an orthogonal matrix $V$ whose columns are right singular vectors.
	\item[(d)] Find an orthogonal matrix $U$ and a $3 \times 2$ diagonal matrix $\Sigma$ such that \[ A = U\Sigma V^T. \]
\end{enumerate}
\subsection{Compute $A^T A$ and find its eigenvalues.}
First, we find the transpose of $A$:
\[
	A^T = \begin{bmatrix} 1 & 1 & 0 \\ 1 & 1 & 0 \end{bmatrix}
\]

Now, compute the product $A^T A$:
\[
	A^T A = \begin{bmatrix} 1 & 1 & 0 \\ 1 & 1 & 0 \end{bmatrix} \begin{bmatrix} 1 & 1 \\ 1 & 1 \\ 0 & 0 \end{bmatrix} = \begin{bmatrix} (1)(1)+(1)(1)+(0)(0) & (1)(1)+(1)(1)+(0)(0) \\ (1)(1)+(1)(1)+(0)(0) & (1)(1)+(1)(1)+(0)(0) \end{bmatrix} = \begin{bmatrix} 2 & 2 \\ 2 & 2 \end{bmatrix}
\]

To find the eigenvalues, we solve the characteristic equation $\det(A^T A - \lambda I) = 0$:
\[
	\det \begin{bmatrix} 2-\lambda & 2 \\ 2 & 2-\lambda \end{bmatrix} = (2-\lambda)^2 - 4 = 0
\]
\[
	4 - 4\lambda + \lambda^2 - 4 = 0
\]
\[
	\lambda^2 - 4\lambda = 0
\]
\[
	\lambda(\lambda - 4) = 0
\]

Therefore, the eigenvalues, ordered from largest to smallest, are:
\[
	\lambda_1 = 4, \quad \lambda_2 = 0
\]

\subsection{Find the singular values of $A$.}
The singular values $\sigma_i$ are the square roots of the eigenvalues of $A^T A$.
\[
	\sigma_1 = \sqrt{\lambda_1} = \sqrt{4} = 2
\]
\[
	\sigma_2 = \sqrt{\lambda_2} = \sqrt{0} = 0
\]

\subsection{Find an orthogonal matrix $V$ whose columns are right singular vectors.}
We need to find the orthonormal eigenvectors for $A^T A$, which will form the columns of $V = [\mathbf{v}_1, \mathbf{v}_2]$.

\textbf{For $\lambda_1 = 4$:}
\[
	(A^T A - 4I)\mathbf{v}_1 = \begin{bmatrix} -2 & 2 \\ 2 & -2 \end{bmatrix} \begin{bmatrix} x \\ y \end{bmatrix} = \begin{bmatrix} 0 \\ 0 \end{bmatrix}
\]
This gives $-2x + 2y = 0 \implies x = y$.
A corresponding unit vector is:
\[
	\mathbf{v}_1 = \begin{bmatrix} 1/\sqrt{2} \\ 1/\sqrt{2} \end{bmatrix}
\]

\textbf{For $\lambda_2 = 0$:}
\[
	(A^T A - 0I)\mathbf{v}_2 = \begin{bmatrix} 2 & 2 \\ 2 & 2 \end{bmatrix} \begin{bmatrix} x \\ y \end{bmatrix} = \begin{bmatrix} 0 \\ 0 \end{bmatrix}
\]
This gives $2x + 2y = 0 \implies x = -y$.
A corresponding unit vector is:
\[
	\mathbf{v}_2 = \begin{bmatrix} -1/\sqrt{2} \\ 1/\sqrt{2} \end{bmatrix}
\]

Placing these as columns, our matrix $V$ is:
\[
	V = \begin{bmatrix} \frac{1}{\sqrt{2}} & -\frac{1}{\sqrt{2}} \\ \frac{1}{\sqrt{2}} & \frac{1}{\sqrt{2}} \end{bmatrix}
\]

\subsection{Find an orthogonal matrix $U$ and a $3 \times 2$ diagonal matrix $\Sigma$ such that $A = U\Sigma V^T$}

\textbf{Find $\Sigma$:}
Since $A$ is a $3 \times 2$ matrix, $\Sigma$ must also be a $3 \times 2$ matrix. We place the singular values on the main diagonal and pad with zeros:
\[
	\Sigma = \begin{bmatrix} \sigma_1 & 0 \\ 0 & \sigma_2 \\ 0 & 0 \end{bmatrix} = \begin{bmatrix} 2 & 0 \\ 0 & 0 \\ 0 & 0 \end{bmatrix}
\]

\textbf{Find $U$:}
The columns of $U = [\mathbf{u}_1, \mathbf{u}_2, \mathbf{u}_3]$ are the left singular vectors. For non-zero singular values, we can find $\mathbf{u}_i$ using the formula $\mathbf{u}_i = \frac{1}{\sigma_i} A \mathbf{v}_i$.

\textbf{For $\sigma_1 = 2$:}
\[
	\mathbf{u}_1 = \frac{1}{2} \begin{bmatrix} 1 & 1 \\ 1 & 1 \\ 0 & 0 \end{bmatrix} \begin{bmatrix} 1/\sqrt{2} \\ 1/\sqrt{2} \end{bmatrix} = \frac{1}{2} \begin{bmatrix} 2/\sqrt{2} \\ 2/\sqrt{2} \\ 0 \end{bmatrix} = \begin{bmatrix} 1/\sqrt{2} \\ 1/\sqrt{2} \\ 0 \end{bmatrix}
\]

Because $\sigma_2 = 0$, we cannot use the formula to find $\mathbf{u}_2$. Instead, we must find vectors $\mathbf{u}_2$ and $\mathbf{u}_3$ such that $\mathbf{u}_1, \mathbf{u}_2, \mathbf{u}_3$ form an orthonormal basis for $\mathbb{R}^3$ (meaning they are unit vectors mutually orthogonal to each other).

To find $\mathbf{u}_2$ orthogonal to $\mathbf{u}_1$, we need a vector $\begin{bmatrix} x \\ y \\ z \end{bmatrix}$ such that $\frac{1}{\sqrt{2}}x + \frac{1}{\sqrt{2}}y + 0z = 0$, which means $x = -y$. Let's pick a normalized vector satisfying this:
\[
	\mathbf{u}_2 = \begin{bmatrix} -1/\sqrt{2} \\ 1/\sqrt{2} \\ 0 \end{bmatrix}
\]

To find $\mathbf{u}_3$ orthogonal to both $\mathbf{u}_1$ and $\mathbf{u}_2$, we can take their cross product, or simply observe that the z-axis unit vector is orthogonal to the xy-plane vectors we just found:
\[
	\mathbf{u}_3 = \begin{bmatrix} 0 \\ 0 \\ 1 \end{bmatrix}
\]

Putting it all together, our orthogonal matrix $U$ is:
\[
	U = \begin{bmatrix} \frac{1}{\sqrt{2}} & -\frac{1}{\sqrt{2}} & 0 \\ \frac{1}{\sqrt{2}} & \frac{1}{\sqrt{2}} & 0 \\ 0 & 0 & 1 \end{bmatrix}
\]

This yields the complete decomposition $A = U\Sigma V^T$.

\end{document}
